\begin{afterword}
\summaryhead

The post opens with Gilbert and Malone's definition: the correspondence bias is inferring lasting dispositions from behaviour that the situation fully explains. When someone else kicks a vending machine we see an angry person; when we kick it, we see the late bus and the lost lunch money.

The author argues from prior probability that the situation is usually the better explanation, since late buses are commoner than people born with extreme anger. Even when people are told that a speaker was assigned a side, they still think the speaker holds it. The author suggests that essences are easier to think of than mechanisms, and mentions false consensus and the fundamental attribution error.

The lesson is to assume that everyone sees themselves as behaving normally, and to ask what situation they believe they are in. The author applies it to people who ask why the author wants to save the world: given the author's beliefs, pressing the green button is the ordinary reaction. Even people who do terrible things are not ``exceptional mutants.''

\responsehead

The post starts from a well-established finding and reports it fairly. It quotes Gilbert and Malone's definition exactly, and its hedged summary of the classic experiment, that people ``don't seem to properly discount'' a situation they have been told about, matches the result: observers inferred attitudes from essays written under instruction, though weaker ones than from essays freely chosen. The advice at the end, to ask what situation people think they are in, follows from that finding.

The claims the post adds to that finding are weaker. The post states as a general rule that we explain our own actions by the situation and others' by disposition. That is the actor-observer hypothesis, which the post runs together with the fundamental attribution error. A meta-analysis of 173 studies, published seven months before the post, found the average difference close to zero. The difference did hold for negative events such as the post's kick, so the post's example fits the evidence, but the general rule does not.

The prior-probability argument compares the situation with an extreme alternative. Late buses are commoner than people born with unnaturally high anger, but ``an angry person'' can mean someone more irritable than most, and such people are common. The comparison does not settle the milder case.

The post also offers its own cause for the bias, that essences are easier to think of than mechanisms, without evidence. Its cited source describes four causes, does not include this one, and says there is no accepted explanation. And the experiment is misdescribed: the essays were about Fidel Castro, not abortion.

The ending applies the lesson to the author's own mission. The post grants that the author's beliefs are unusual and call for explanation, and answers the worry about false consensus with a bet. The post argues only that the reaction is normal given the beliefs. It leaves the case for the beliefs to the chapter cited in its footnote.

In short: a real finding, reported fairly, surrounded by a self-other rule that the evidence of the time did not support in general, a prior-probability argument against an extreme alternative, and an unsupported cause.
\end{afterword}
